\documentclass[11pt]{article}
\usepackage[a4paper,margin=25mm]{geometry}
\usepackage{amsmath,amssymb,mathtools}
\usepackage{fontspec,unicode-math}
\setmainfont{Latin Modern Roman}
\setsansfont{Latin Modern Sans}
\setmonofont{Latin Modern Mono}
\setmathfont{Latin Modern Math}
\usepackage{xurl,hyperref}
\hypersetup{colorlinks=true,urlcolor=blue,linkcolor=blue}
\newcommand{\sref}[2]{\hyperlink{source:#1:#2}{[S#2]}}
\newcommand{\cataloguescope}[1]{\paragraph{Mathematical question.}#1}
\begin{document}
\section{Conjecture 4.16 — Quantitative Refinement of Schanuel's Conjecture}\label{statement}

\subsection{Definitions and mathematical statement}

\paragraph{Shared notation.}$\mathbb N=\{1,2,\ldots\}$, and $\mathbb Z,\mathbb Q,\mathbb R,\mathbb C$ denote the integers, rationals, reals and complex numbers. Any other conventions are specified locally.


Let $n,d\ge1$. For an integer polynomial $P$ in $2n$ variables, its degree means total degree and its usual height is the largest absolute coefficient. In $R=\mathbb Q[X_1,\ldots,X_n,Y_1,\ldots,Y_n]$, the rank of a prime ideal is its height (the supremum of lengths of strict chains of prime ideals ending at it); the rank of an ideal $I$ is $\min_{\mathfrak p\supset I}\operatorname{ht}\mathfrak p$.
\paragraph{Statement recorded in the supplied source.}
Let $x_1,\ldots,x_n\in\mathbb C$ be linearly independent over $\mathbb Q$ and suppose
\[
\forall\varepsilon>0\ \exists H_0>0\ \forall H\ge H_0\ \forall(h_1,\ldots,h_n)\in\mathbb Z^n,
\quad 0<\max_i|h_i|\le H\Longrightarrow
\left|\sum_i h_ix_i\right|\ge e^{-H^{\varepsilon}}.
\]
Is there a constant $C=C(x_1,\ldots,x_n,d)>0$ such that, for every integer $H\ge2$ and every $P_1,\ldots,P_{n+1}\in\mathbb Z[X_1,\ldots,X_n,Y_1,\ldots,Y_n]$ of total degrees at most $d$ and heights at most $H$ whose ideal in $R$ has rank $n+1$,
\[
\sum_{j=1}^{n+1}|P_j(x_1,\ldots,x_n,e^{x_1},\ldots,e^{x_n})|\ge H^{-C}?
\] \sref{statement}{2}

\subsection{Short English statement}
Under the stated quantitative independence hypothesis, must polynomial systems of the specified codimension stay at least polynomially far from vanishing on the exponential tuple?
\subsection{Sources}
\begin{description}
\item[\hypertarget{source:statement:1}{S1}] ULAM.AI and the UnsolvedMath Contributors, UnsolvedMath: A Curated Collection of Open Mathematics Problems, catalogue entry. CC BY 4.0. \url{https://huggingface.co/datasets/ulamai/UnsolvedMath}.
\item[\hypertarget{source:statement:2}{S2}] Michel Waldschmidt, \emph{Open Diophantine Problems}, arXiv:math/0312440v2 (24 January 2004), Conjecture 4.16, p.42; ideal-rank definition, p.41. \url{https://arxiv.org/abs/math/0312440v2}.
\end{description}
\label{end:statement}
\end{document}
